\documentclass{article}
\usepackage[margin=1.5cm]{geometry}
\usepackage{amsmath}
\usepackage{enumitem}

\setlist{nosep,leftmargin=*}

\begin{document}
\twocolumn

\title{Chapter 10 Warm-Up: List and Dictionary Comprehensions}
\author{Prof. Jordan C. Hanson}
\date{}

\maketitle

\section{A List of Student Names}

Suppose we are given a list of student names:

\begin{verbatim}
names = ["Ana","Ben","Chloe","Diego","Elena"]
\end{verbatim}
\noindent
Student IDs begin with the eight-digit number \texttt{02601001} and are assigned sequentially.  Thus, the first student receives ID \texttt{02601001}, the second receives \texttt{02601002}, etc.

\section{Generating IDs with a Loop}

One way to create the student IDs is

\begin{verbatim}
st_ids = []
first = 02601001
for i in range(len(names)):
    st_ids.append(first + i)
\end{verbatim}

\begin{enumerate}
\item What values does \texttt{i} take?
\item Predict the contents of \texttt{student\_ids}.
\item Why is \texttt{len(names)} useful here?
\end{enumerate}

\section{List Comprehension}

Python can express the same operation more compactly:

\begin{verbatim}
st_ids = [first+i for i in range(len(names))]
\end{verbatim}
\noindent
This is called a \textbf{list comprehension}, with the pattern

\begin{verbatim}
[expression for item in sequence]
\end{verbatim}
\noindent
For each value of \texttt{i}, Python evaluates \texttt{first+i} and places the result into the list.

\section{Building a Student Dictionary}

Now suppose we want a dictionary in which

\begin{itemize}
\item the \textbf{student ID} is the key, and
\item the \textbf{student name} is the value.
\end{itemize}
\noindent
Using an ordinary loop:

\begin{verbatim}
students = {}
for i, name in enumerate(names):
    students[first+i] = name
\end{verbatim}
\noindent
The dictionary will contain entries such as

\begin{verbatim}
{
 2601001: "Ana",
 2601002: "Ben",
 2601003: "Chloe",
 ...
}
\end{verbatim}
\noindent
The function \texttt{enumerate(names)} provides both the index and the name for each student.

\section{Dictionary Comprehension}

The entire dictionary can instead be created with

\begin{verbatim}
students = {
    first+i: name
    for i, name in enumerate(names)
}
\end{verbatim}
\noindent
This is a \textbf{dictionary comprehension}, with the pattern

\begin{verbatim}
{
 key_expression: value_expression
 for item in sequence
}
\end{verbatim}

\section{Compare the Two}

A list comprehension:

\begin{verbatim}
[first_id + i
 for i in range(len(names))]
\end{verbatim}
\noindent
creates a sequence of student IDs.  A dictionary comprehension creates the ID-name pairs:

\begin{verbatim}
{first_id + i: name
 for i, name in enumerate(names)}
\end{verbatim}

\section{Check Your Understanding}

\begin{enumerate}
\item Which comprehension creates only the student IDs?

\item Which comprehension associates each ID with a name?

\item Where do the names in the dictionary come from?

\item Which part of the dictionary comprehension generates
the ID?

\item What would the ID for \texttt{"Diego"} be?

\item If a sixth name is added to \texttt{names}, does either
comprehension need to be rewritten? Explain.
\end{enumerate}

\end{document}
